
\documentclass{article}
\usepackage{colortbl}
\usepackage{makecell}
\usepackage{multirow}
\usepackage{supertabular}

\begin{document}

\newcounter{utterance}

\centering \large Interaction Transcript for game `matchit\_ascii', experiment `different\_grid', episode 3 with NEW\_Reflect\_Agent.
\vspace{24pt}

{ \footnotesize  \setcounter{utterance}{1}
\setlength{\tabcolsep}{0pt}
\begin{supertabular}{c@{$\;$}|p{.15\linewidth}@{}p{.15\linewidth}p{.15\linewidth}p{.15\linewidth}p{.15\linewidth}p{.15\linewidth}}
    \# & $\;$A & \multicolumn{4}{c}{Game Master} & $\;\:$B\\
    \hline

    \theutterance \stepcounter{utterance}  
    & & \multicolumn{4}{p{0.6\linewidth}}{
        \cellcolor[rgb]{0.9,0.9,0.9}{
            \makecell[{{p{\linewidth}}}]{
                \texttt{\tiny{[P1$\langle$GM]}}
                \texttt{You are participating in a collaborative guessing game.} \\
\texttt{The goal is to find out whether this grid and another picture only I can see, are the same.} \\
\texttt{Please describe your grid first. Then, I will provide my description and we can ask each other questions about the grids to figure out whether they are the same. Now start your short grid description with "DESCRIPTION:" followed by the description. Do not add anything else. This is your grid:} \\
\texttt{X ▢ ▢ ▢ X} \\
\texttt{X ▢ ▢ ▢ X} \\
\texttt{X ▢ X ▢ X} \\
\texttt{X X ▢ X X} \\
\texttt{X ▢ ▢ ▢ X} \\
            }
        }
    }
    & \\ \\

    \theutterance \stepcounter{utterance}  
    & \multicolumn{4}{p{0.45\linewidth}}{
        \cellcolor[rgb]{0.95,0.95,0.95}{
            \makecell[{{p{\linewidth}}}]{
                \texttt{\tiny{[P1$|$internal]}}
                \texttt{Reflections from previous episodes:} \\
\texttt{{-}} \\
\texttt{ Failed. I gave a clear description of my grid (border of X's with interior spaces). The user then described a different pattern (central cross shape). Our grids were not the same, which we determined after exchanging descriptions. My initial description was accurate and sufficient for comparison. Next time, I would follow the same approach: provide a concise, accurate initial description to enable quick identification of match or mismatch.} \\
\texttt{{-}} \\
\texttt{ Failed. I described the grid but did not ask clarifying questions about the other player's hidden grid to compare effectively. Next time, after giving my description, I should actively ask specific questions about their grid's layout to identify matches or mismatches.} \\
\texttt{{-}} \\
\texttt{ Failed. I described the grid as having a central row with three X's and two squares, but the user's unseen grid likely differed. I should have asked clarifying questions about the user's grid to compare details, rather than only providing my own description. Next time, I will immediately follow my description with specific questions about their grid's pattern to identify mismatches faster.} \\
            }
        }
    }
    & & \\ \\

    \theutterance \stepcounter{utterance}  
    & \multicolumn{4}{p{0.45\linewidth}}{
        \cellcolor[rgb]{0.95,0.95,0.95}{
            \makecell[{{p{\linewidth}}}]{
                \texttt{\tiny{[P1$|$internal]}}
                \texttt{{[}attempt 1 – valid{]}} \\
\texttt{DESCRIPTION: A 5x5 grid with X's forming a border and a central cross pattern. The outer rows and columns are mostly X's, with the center cell (3,3) being X, and the middle row (3,1), (3,2), (3,4), (3,5) are X's, while (3,3) is also X. The interior has X's in the center row and center column, forming a cross, with the corners and edges filled with X's. Specifically:  } \\
\texttt{Row 1: X ▢ ▢ ▢ X  } \\
\texttt{Row 2: X ▢ ▢ ▢ X  } \\
\texttt{Row 3: X ▢ X ▢ X  } \\
\texttt{Row 4: X X ▢ X X  } \\
\texttt{Row 5: X ▢ ▢ ▢ X} \\
            }
        }
    }
    & & \\ \\

    \theutterance \stepcounter{utterance}  
    & \multicolumn{4}{p{0.6\linewidth}}{
        \cellcolor[rgb]{0.9,0.9,0.9}{
            \makecell[{{p{\linewidth}}}]{
                \texttt{\tiny{[P1$\rangle$GM]}}
                \texttt{DESCRIPTION: A 5x5 grid with X's forming a border and a central cross pattern. The outer rows and columns are mostly X's, with the center cell (3,3) being X, and the middle row (3,1), (3,2), (3,4), (3,5) are X's, while (3,3) is also X. The interior has X's in the center row and center column, forming a cross, with the corners and edges filled with X's. Specifically:  } \\
\texttt{Row 1: X ▢ ▢ ▢ X  } \\
\texttt{Row 2: X ▢ ▢ ▢ X  } \\
\texttt{Row 3: X ▢ X ▢ X  } \\
\texttt{Row 4: X X ▢ X X  } \\
\texttt{Row 5: X ▢ ▢ ▢ X} \\
            }
        }
    }
    & & \\ \\

    \theutterance \stepcounter{utterance}  
    & & & \multicolumn{2}{p{0.3\linewidth}}{
        \cellcolor[rgb]{0.9,0.9,0.9}{
            \makecell[{{p{\linewidth}}}]{
                \texttt{\tiny{[GM$|$GM]}}
                \texttt{continue} \\
            }
        }
    }
    & & \\ \\

    \theutterance \stepcounter{utterance}  
    & & \multicolumn{4}{p{0.6\linewidth}}{
        \cellcolor[rgb]{0.9,0.9,0.9}{
            \makecell[{{p{\linewidth}}}]{
                \texttt{\tiny{[GM$\rangle$P2]}}
                \texttt{You are participating in a collaborative guessing game.} \\
\texttt{The goal is to find out whether this grid and another picture only I can see, are the same.} \\
\texttt{Please describe your grid first. Then, I will provide my description and we can ask each other questions about the grids to figure out whether they are the same. Now start your short grid description with "DESCRIPTION:" followed by the description. Do not add anything else. This is your grid:} \\
\texttt{▢ ▢ X ▢ ▢} \\
\texttt{▢ ▢ X ▢ ▢} \\
\texttt{X X X X X} \\
\texttt{▢ ▢ X ▢ ▢} \\
\texttt{▢ ▢ X ▢ ▢} \\
            }
        }
    }
    & \\ \\

    \theutterance \stepcounter{utterance}  
    & & & \multicolumn{4}{p{0.45\linewidth}}{
        \cellcolor[rgb]{0.95,0.95,0.95}{
            \makecell[{{p{\linewidth}}}]{
                \texttt{\tiny{[P2$|$internal]}}
                \texttt{Reflections from previous episodes:} \\
\texttt{{-}} \\
\texttt{ Failed. I described a diamond shape, but user described a cross shape. I should have asked about the pattern earlier instead of just describing. Next time, ask a clarifying question first, like "QUESTION: Is the pattern symmetrical?" to quickly identify major differences.} \\
\texttt{{-}} \\
\texttt{ FAIL. My description was wrong: I described the grid as having vertical X stripes, but it actually had horizontal X stripes. This mismatch likely confused the user. Next time, I must carefully observe the grid orientation before describing.} \\
\texttt{{-}} \\
\texttt{ Failed. My description incorrectly stated the grid had X's in all positions except row 3, which was wrong. The actual grid was a right triangle of X's starting from the top{-}left. My error misled the partner from the start. Next time, I must carefully observe the pattern before describing.} \\
            }
        }
    }
     \\ \\

    \theutterance \stepcounter{utterance}  
    & & & \multicolumn{4}{p{0.45\linewidth}}{
        \cellcolor[rgb]{0.95,0.95,0.95}{
            \makecell[{{p{\linewidth}}}]{
                \texttt{\tiny{[P2$|$internal]}}
                \texttt{{[}attempt 1 – valid{]}} \\
\texttt{DESCRIPTION: The grid has a cross shape, with X's forming a vertical line in the middle column and a horizontal line in the middle row, intersecting at the center. All other positions are empty (□).} \\
            }
        }
    }
     \\ \\

    \theutterance \stepcounter{utterance}  
    & & & \multicolumn{4}{p{0.6\linewidth}}{
        \cellcolor[rgb]{0.9,0.9,0.9}{
            \makecell[{{p{\linewidth}}}]{
                \texttt{\tiny{[GM$\langle$P2]}}
                \texttt{DESCRIPTION: The grid has a cross shape, with X's forming a vertical line in the middle column and a horizontal line in the middle row, intersecting at the center. All other positions are empty (□).} \\
            }
        }
    }
     \\ \\

    \theutterance \stepcounter{utterance}  
    & & & \multicolumn{2}{p{0.3\linewidth}}{
        \cellcolor[rgb]{0.9,0.9,0.9}{
            \makecell[{{p{\linewidth}}}]{
                \texttt{\tiny{[GM$|$GM]}}
                \texttt{continue} \\
            }
        }
    }
    & & \\ \\

    \theutterance \stepcounter{utterance}  
    & & \multicolumn{4}{p{0.6\linewidth}}{
        \cellcolor[rgb]{0.9,0.9,0.9}{
            \makecell[{{p{\linewidth}}}]{
                \texttt{\tiny{[GM$\rangle$P2]}}
                \texttt{This is my DESCRIPTION: A 5x5 grid with X's forming a border and a central cross pattern. The outer rows and columns are mostly X's, with the center cell (3,3) being X, and the middle row (3,1), (3,2), (3,4), (3,5) are X's, while (3,3) is also X. The interior has X's in the center row and center column, forming a cross, with the corners and edges filled with X's. Specifically:  } \\
\texttt{Row 1: X ▢ ▢ ▢ X  } \\
\texttt{Row 2: X ▢ ▢ ▢ X  } \\
\texttt{Row 3: X ▢ X ▢ X  } \\
\texttt{Row 4: X X ▢ X X  } \\
\texttt{Row 5: X ▢ ▢ ▢ X} \\
\texttt{Now ask a question in order to find out new aspects of my grid that may be different to your grid. Start with "QUESTION:" and do not add anything else.} \\
            }
        }
    }
    & \\ \\

    \theutterance \stepcounter{utterance}  
    & & & \multicolumn{4}{p{0.6\linewidth}}{
        \cellcolor[rgb]{0.9,0.9,0.9}{
            \makecell[{{p{\linewidth}}}]{
                \texttt{\tiny{[GM$\langle$P2]}}
                \texttt{DESCRIPTION: ▢ ▢ X ▢ ▢  } \\
\texttt{▢ ▢ X ▢ ▢  } \\
\texttt{X X X X X  } \\
\texttt{▢ ▢ X ▢ ▢  } \\
\texttt{▢ ▢ X ▢ ▢  } \\
\\ 
\texttt{QUESTION: In your grid, are the corners (1,1), (1,5), (5,1), and (5,5) all X's?} \\
            }
        }
    }
     \\ \\

    \theutterance \stepcounter{utterance}  
    & & & \multicolumn{2}{p{0.3\linewidth}}{
        \cellcolor[rgb]{0.9,0.9,0.9}{
            \makecell[{{p{\linewidth}}}]{
                \texttt{\tiny{[GM$|$GM]}}
                \texttt{abort, first word: DESCRIPTION:, but expected QUESTION:} \\
            }
        }
    }
    & & \\ \\

\end{supertabular}
}

\end{document}
